\begin{zhabstract}
  摘要二字间间距为两个字符，\strong{中文摘要}正文，每段开头\strong{左起空两个字符}，
  段与段之间不空行，\strong{硕士论文中文摘要}正文字数 600-800。
  \strong{由申请人设计的关键词}，须在该关键词的\strong{右上角标注*}，
  并在该页的页脚注明一句话：“*表示非汉语主题词表用词”。

  摘要正文内容一般包括：从事这项研究工作的目的和意义；作者独立进行的研究工作的概括性叙述；
  研究获得的主要结论或提出的主要观点。硕士学位论文摘要应突出论文的新见解，
  博士学位论文摘要应突出论文的创新点。(宋体，小四，1.2 倍行距)
  \end{zhabstract}

\zhkeywords{\strong{关键词由 3\~{}5 个组成}；XXXX；XXXX；XXXX}
\zhtypedisseration{XXXXX （下列类型十选一）}

\begin{enabstract}
  英文每段开头\strong{左对齐顶格}编排，段与段之间不空行，
  (Times New Roma，小四，1.2 倍行距)。
  \strong{由申请人设计的关键词}，须在该关键词的\strong{右上角标注*}，
  并在该页的\strong{页脚注明一句话}：“*表示非汉语主题词表用词”。

  English paragraph at the beginning of the left aligned top layout,
  paragraph and paragraph between the no empty line,
  (Time New Roma, minor four, 1.2 times the space).
  Keywords designed by the applicant should be marked with * in the upper
  right corner of the keyword and a sentence in the footer of the page:
  "* indicates non-Chinese thesaurus words".
\end{enabstract}

\enkeywords{ XXX; XXX; XXX; XXX (Times New Roma，五号，加黑)}
\entypedisseration{XXXXX (下列类型十选一)}
